\originalsection{点变换不改变方程的形式}
同一物理系统可用不同广义坐标. 设$Q=Q(q,t)$是局部可逆点变换, 定义$L'(Q,\dot Q,t)=L(q,\dot q,t)$. 固定时刻$\delta Q_j=Q_{j,q_i}\delta q_i$, 在作用量的一阶变分中比较任意$\delta q_i$, 得
\[
\frac{\dd}{\dd t}L_{\dot q_i}-L_{q_i}
=\sum_j Q_{j,q_i}\left(\frac{\dd}{\dd t}L'_{\dot Q_j}-L'_{Q_j}\right).
\]
Jacobi矩阵可逆时, 一组Euler--Lagrange式为零等价于另一组为零. 时间相关的点变换也成立, 但新拉格朗日量须保留$\dot Q=Q_q\dot q+Q_t$带来的项. 极坐标就是例子: $L=m(\dot r^2+r^2\dot\phi^2)/2-U(r,\phi)$给$m\ddot r-mr\dot\phi^2=-U_r$与$\dd(mr^2\dot\phi)/\dd t=-U_\phi$. 后式右端是转矩, 不是沿圆周的力; 圆周力为$-U_\phi/r$. 换坐标与改用旋转参考系是不同选择, 应分别核查.

\originalsection{圆锥与可移动支点}
圆锥半顶角为$0<\alpha<\pi/2$, 轴竖直向上, 表面$z=r\cot\alpha$. 光滑表面上的质点用$(r,\phi)$即可描述, 消去$z$得
\[
L=\frac m2\left(\csc^2\alpha\,\dot r^2+r^2\dot\phi^2\right)
-mgr\cot\alpha.
\]
角动量$J=mr^2\dot\phi$守恒, 径向方程为$m\csc^2\alpha\,\ddot r=J^2/(mr^3)-mg\cot\alpha$. 圆轨道半径满足$r_0^3=J^2/(m^2g\cot\alpha)$; 在$r_0$附近线性化得$\omega_r^2=3g\sin\alpha\cos\alpha/r_0>0$. 计算只消去了反力, 没有自动证明质点仍接触表面; 有单侧接触时须另验法向力非负.

再取质量$M$的小车在水平光滑轨上移动, 质量$m$的摆锤以长$\ell$的杆连到小车, $\theta$从竖直向下量. 质点坐标为$(x+\ell\sin\theta,-\ell\cos\theta)$, 所以
\[
L=\frac{M+m}{2}\dot x^2+m\ell\cos\theta\,\dot x\dot\theta
+\frac{m\ell^2}{2}\dot\theta^2+mg\ell\cos\theta.
\]
交叉项来自摆锤同时跟随小车和绕支点运动. 方程为$(M+m)\ddot x+m\ell(\cos\theta\ddot\theta-\sin\theta\dot\theta^2)=0$及$\ell\ddot\theta+\cos\theta\ddot x+g\sin\theta=0$. 第一式积分给总水平动量守恒, 不意味着小车速度恒定. 在线性小振幅且总动量为零的情形, $\ddot x=-m\ell\ddot\theta/(M+m)$, 从而$\ddot\theta+[(M+m)g/(M\ell)]\theta=0$. 自由支点会使频率高于固定支点的$\sqrt{g/\ell}$; 第6章将用这个质量矩阵演示Hamilton化.

\originalsection{滚动何时可以积分成位置约束}
半径$a$、转动惯量$I=\kappa ma^2$的轮沿固定斜面直线滚动. 以$x$沿坡向下、$\varphi$沿实际转动方向, 无滑动给$\dot x=a\dot\varphi$, 可积分成$x-a\varphi=$常数. 消元后$T=m(1+\kappa)\dot x^2/2$, 得加速度$\ddot x=g\sin\beta/(1+\kappa)$. 平动式$m\ddot x=mg\sin\beta-f$和转动式$I\ddot\varphi=fa$给静摩擦向上, $f=mg\kappa\sin\beta/(1+\kappa)$. 所需条件是$\mu_s\ge\kappa\tan\beta/(1+\kappa)$; 条件不满足就必须改为滑动模型. 静摩擦虽非零, 理想固定接触点瞬时静止, 不作总功.

若竖直轮可在平面上改变朝向$\theta$, 以该朝向定义前进单位矢量$(\cos\theta,\sin\theta)$, 速度约束变为$\dot x=a\cos\theta\dot\varphi$, $\dot y=a\sin\theta\dot\varphi$. 允许速度方向张成向量场$X=\partial_\theta$与$Y=\partial_\varphi+a\cos\theta\partial_x+a\sin\theta\partial_y$. 它们的交换子为$[X,Y]=-a\sin\theta\partial_x+a\cos\theta\partial_y$. 若它在$X,Y$的张成空间内, 无$\theta,\varphi$分量迫使两个系数都为零, 但其水平分量不为零, 矛盾. 因而允许切平面分布不可积为一个位置约束子流形. 此为中文补充的局部可积性检查; 不能只因某个约束写在速度上就判断非完整, 直线滚动已是反例.

\originalsection{单侧接触与摩擦耗散}
单侧接触写成间隙$g(q)\ge0$, 法向反力$N\ge0$, 并满足$Ng(q)=0$. 接触阶段用$g=0$的约束方程计算$N$; 达到零后须检查是否进入分离阶段. 例如质点从光滑固定球顶受微扰向下滑, 从竖直量角$\theta$, 能量给$v^2=2ga(1-\cos\theta)$, 径向式给$N=mg\cos\theta-mv^2/a=mg(3\cos\theta-2)$. 到$\cos\theta=2/3$离开球面. 此后不能继续沿球面解并接受负反力. 静止于精确球顶的理想初值是平衡解, 所说下滑是受扰动或极限释放状态.

速度线性阻力$\vect F_a=-b_a\vect v_a$可用Rayleigh函数$\mathcal R=\sum_ab_av_a^2/2$写广义力$Q_i=-\mathcal R_{\dot q_i}$. 这是利用$\partial\vect v_a/\partial\dot q_i=\partial\vect r_a/\partial q_i$, 将矢量摩擦化为标量导数. 对固定约束且$L$不显含时间,$\dot E=\sum_iQ_i\dot q_i=-2\mathcal R$; 移动约束时还要保留外部输入功. 更一般的各向同性阻力$-h_a(v_a)\vect v_a/v_a$对应$\mathcal R=\sum_a\int_0^{v_a}h_a(s)\dd s$, 但其能量损失不一定等于$2\mathcal R$. 干摩擦的$h=\mu N$使函数在零速度处不可微, 静止时应使用$|F_s|\le\mu_sN$的静摩擦条件, 不能把符号函数在零点随意设零就完成动力学.

球在黏性流体中缓慢沉降是一个例子. 向下速度为$v$, 球质量$m$, 排开流体质量$m_f$, 准稳态Stokes阻力近似下
\[
m\dot v=(m-m_f)g-6\pi\eta av,
\qquad v(t)=v_\infty+[v(0)-v_\infty]e^{-t/\tau},
\]
其中$v_\infty=(m-m_f)g/(6\pi\eta a)$, $\tau=m/(6\pi\eta a)$. 浮力改变驱动力, 并不把球的惯性质量替成$m-m_f$; 原稿第21页如此替换须纠正. 若非定常流体附加质量和历史力不可忽略, 则上述简单模型仍须扩展. Stokes阻力本身将在第9章从流体边值问题推出.
